%!TEX TS-program = xelatex
%!TEX encoding = UTF-8 Unicode
%%%%%%%%%%%%%%%%%%%%%%%%%%%%%%%%%%%%%%%%%%%%%%%%%%%%%%%%%%%%
%% VIRTIO 1.0 Specification Draft
%%%%%%%%%%%%%%%%%%%%%%%%%%%%%%%%%%%%%%%%%%%%%%%%%%%%%%%%%%%%

\documentclass[10pt,titlepage]{report}

\usepackage[top=1in,left=1in,bottom=0.5in,right=1in,paperwidth=8.5in, paperheight=11in]{geometry}
\usepackage{parskip}
\usepackage{setspace}
\usepackage[english]{babel}
\usepackage{ae,aecompl} 
\usepackage[T1]{fontenc}
\usepackage[xetex]{graphicx}
\usepackage{placeins}
\usepackage{framed}
\usepackage[xetex]{hyperref}
\usepackage{lineno} 
\usepackage{listings}
\usepackage{courier}
\usepackage{fancyvrb}
\usepackage{array}
\usepackage{longtable}
\usepackage{multirow}
\usepackage{rotating}
\usepackage{fancyhdr}
\usepackage{lastpage}
\usepackage{todonotes}
\usepackage[explicit]{titlesec}
\usepackage{mdwlist}
\usepackage{url}
\usepackage{xifthen}
\usepackage[no-math]{fontspec}
\usepackage{xltxtra}
\usepackage{etoolbox}
\usepackage{tabularx}
\usepackage{underscore}
\usepackage{xstring}
\usepackage{ellipsis}
% This is an alternative to package underscore above
% It makes it possible to disable hyphenation of upper-case
% identifiers with underscore (of which we have many) by
% treating them as acronyms.
% However, the result doesn't look pretty, so don't use this
% for now.
%\catcode`\_=\active
%\def_{\_}
% Don't hyphenate acronyms
%\uchyph=0
\setmainfont[Mapping=tex-text]{Arial}
\setromanfont{Arial}
\setmonofont{Courier New}

\urlstyle{rm}

% commands
% define VIRTIO Working Draft number and date
\newcommand{\virtioversion}{Virtual I/O Device (VIRTIO) Version 1.0}
\newcommand{\virtioworkingdraft}{01}
\newcommand{\virtioworkingdraftdate}{05 Nov 2013}
\newcommand{\virtiospecfile}{virtio-v1.0-wd01}
\newcommand{\virtioworkproduct}{Standards Track Draft}
\newcommand{\virtiooasiscopyright}{Copyright © OASIS Open 2013. All Rights Reserved.}

% PDF setup
\hypersetup{
pdfauthor = {OASIS Virtual I/O Device (VIRTIO) TC},
pdftitle = {Virtual I/O Device (VIRTIO) Version 1.0},
pdfsubject = {Virtual I/O Device (VIRTIO) Version 1.0},
pdfkeywords = {VIRTIO, Virtual I/O Device},
colorlinks = {true},
linkcolor=blue,
urlcolor = blue}

%%%%%%%%%%%%%%%%%%%%%%%%%%%%%%%%%%%%%%%%%%%%%%%%%%%%%%%%%%%%
%% Commands for the VIRTIO specification document
%%%%%%%%%%%%%%%%%%%%%%%%%%%%%%%%%%%%%%%%%%%%%%%%%%%%%%%%%%%%

\definecolor{oasis1}{RGB}{85,38,129}
\definecolor{oasis2}{RGB}{227,175,27}
\definecolor{shadecolor}{RGB}{230,230,230}

% How we format a field name
\newcommand{\field}[1]{\emph{#1}}

% Mark a normative paragraph (driver or device)
\newcommand{\drivernormative}[1]{\phantomsection\label{drivernormative:#1}}
\newcommand{\devicenormative}[1]{\phantomsection\label{devicenormative:#1}}

%%%%%%%%%%%%%%%%%%%%%%%%%%%%%%%%%%%%%%%%%%%%%%%%%%%%%%%%%%%%
%% PDF specfic commands
%%%%%%%%%%%%%%%%%%%%%%%%%%%%%%%%%%%%%%%%%%%%%%%%%%%%%%%%%%%%

\newcommand{\virtiographics}[3]
{
\begin{figure}[htb]
\centering
\includegraphics[#2]{images/#1.pdf}
\caption{#3}
\end{figure}
\FloatBarrier
}


%%%%%%%%%%%%%%%%%%%%%%%%%%%%%%%%%%%%%%%%%%%%%%%%%%%%%%%%%%%%
%%% VIRTIO title sections
%%%%%%%%%%%%%%%%%%%%%%%%%%%%%%%%%%%%%%%%%%%%%%%%%%%%%%%%%%%%

\newcommand{\oasistitle}[1]
{
\begin{spacing}{1.5}
\fontspec[Color=552681]{Arial}
\fontsize{24}{24}
\textbf{\noindent#1}
\end{spacing}
}

\newcommand{\oasisstagelabel}[1]
{
\begin{spacing}{1.1}
\fontspec[Color=552681]{Arial}
\fontsize{18}{18}
\textbf{\noindent#1}
\end{spacing}
}

\newcommand{\oasisdate}[1]
{
\begin{spacing}{1.1}
\fontspec[Color=552681]{Arial}
\fontsize{18}{18}
\textbf{\noindent#1}
\end{spacing}
}

\newcommand{\oasisnoticelabel}[1]
{
\begin{spacing}{1.5}
\fontspec[Color=552681]{Arial}
\fontsize{18}{18}
\textbf{\noindent#1}
\end{spacing}
}

\newcommand{\oasisspecificationuris}[1]
{
{
\fontspec[Color=552681]{Arial}
\fontsize{12}{12}
\textbf{\noindent#1}
}
}

\newenvironment{oasistitlesection}[1]
{
\begin{tabular}{p{0in}p{5.8in}}
\multicolumn{2}{l}{\fontspec[Color=552681]{Arial}\fontsize{10}{10}\textbf{#1:}}\\[-0.05in]
&\setlength{\topsep}{0in}\setlength{\parsep}{0in}\setlength{\partopsep}{0in}   }
{
\end{tabular}
\vspace{-0.11in}
}



%%%%%%%%%%%%%%%%%%%%%%%%%%%%%%%%%%%%%%%%%%%%%%%%%%%%%%%%%%%%
%%% VIRTIO type attributes
%%%%%%%%%%%%%%%%%%%%%%%%%%%%%%%%%%%%%%%%%%%%%%%%%%%%%%%%%%%%

% new environment: VIRTIO type attribute
\newenvironment{virtioattribute}[2]
{
\vspace{0.05in}
\begin{tabular}{p{2in}l}
\texttt{\textbf{#1}} & #2\\
\end{tabular}\\\
\begin{tabular}{p{0.3in}p{5.7in}}
&}
{\\
\end{tabular}\\
\vspace{0.05in}
}

% new command: VIRTIO type attribute value
\newcommand{\virtioattributevalue}[2]
{
\vspace{0.05in}
\noindent
\texttt{\textbf{#1}}\\
\indent
Value: #2
\vspace{0.05in}
}


%%%%%%%%%%%%%%%%%%%%%%%%%%%%%%%%%%%%%%%%%%%%%%%%%%%%%%%%%%%%
%%% VIRTIO properties
%%%%%%%%%%%%%%%%%%%%%%%%%%%%%%%%%%%%%%%%%%%%%%%%%%%%%%%%%%%%

% new environment: VIRTIO propery
\newenvironment{virtioproperty}[2]
{
\vspace{0.05in}
\begin{tabular}{p{0.1in}p{2.3in}p{3.5in}}
\multicolumn{2}{p{1.8in}}{\textbf{\texttt{#1}}} & #2\vspace{0.05in}\\
}
{
\end{tabular}\\
\vspace{0.05in}\\
}

% new command: VIRTIO property detail
\newcommand{\virtiopropertydetail}[2]{& #1: & #2\\}

% new command: VIRTIO property comment
\newcommand{\virtiopropertycomment}[1]{\\&\multicolumn{2}{p{5.7in}}{#1}\\}


%%%%%%%%%%%%%%%%%%%%%%%%%%%%%%%%%%%%%%%%%%%%%%%%%%%%%%%%%%%%
%%% VIRTIO allowable actions
%%%%%%%%%%%%%%%%%%%%%%%%%%%%%%%%%%%%%%%%%%%%%%%%%%%%%%%%%%%%

% new environment: VIRTIO allowable action
\newenvironment{virtioallowableaction}[1]
{
\vspace{0.05in}
\noindent
\textbf{#1}
\vspace{0.05in}\\
\begin{tabular}{p{0.1in}p{1in}p{4.8in}}
}
{
\end{tabular}\\
\vspace{0.05in}\\
}

% new command: VIRTIO allowable action detail
\newcommand{\virtioallowableactiondetail}[2]{& #1: & #2\\}


%%%%%%%%%%%%%%%%%%%%%%%%%%%%%%%%%%%%%%%%%%%%%%%%%%%%%%%%%%%%
%%% VIRTIO version properties
%%%%%%%%%%%%%%%%%%%%%%%%%%%%%%%%%%%%%%%%%%%%%%%%%%%%%%%%%%%%

% new environment: VIRTIO version property
\newenvironment{virtioversionproperty}[2]
{
\vspace{0.05in}
\begin{tabular}{p{3in}l}
\texttt{\textbf{#1}} & #2\\
\end{tabular}\\\
\begin{tabular}{p{0.3in}p{5.7in}}
&}
{\\
\end{tabular}\\
\vspace{0.05in}
}

%%%%%%%%%%%%%%%%%%%%%%%%%%%%%%%%%%%%%%%%%%%%%%%%%%%%%%%%%%%%
%%% VIRTIO general description
%%%%%%%%%%%%%%%%%%%%%%%%%%%%%%%%%%%%%%%%%%%%%%%%%%%%%%%%%%%%

% new environment: VIRTIO description
\newenvironment{virtiodescription}[1]
{
\vspace{0.05in}
\noindent
\begin{tabular}{p{0.1in}p{5.9in}}
\multicolumn{2}{l}{\textbf{#1:}}\\
&\setlength{\topsep}{0in}\setlength{\parsep}{0in}\setlength{\partopsep}{0in}   }
{\\
\end{tabular}\\
\vspace{0.05in}
}


\newenvironment{virtioexample}[1]
{
\vspace{0.2in}
\begin{shaded}
\textbf{Example:}\\
Request: \texttt{#1-request.log}\\
Response: \texttt{#1-response.log}
\end{shaded}
\vspace{0.2in}
}


% listings setup
\definecolor{LISTING_background}{RGB}{245,245,245}
\lstset{
	basicstyle=\footnotesize\ttfamily,
	extendedchars=true,
	backgroundcolor=\color{LISTING_background},
	showspaces=false, 
	showstringspaces=false,
	showtabs=false, 
	frame=single, 
	tabsize=2,  
	breaklines=true, 
	breakatwhitespace=true,
	postbreak=\raisebox{0ex}[0ex][0ex]{\ensuremath{\hookrightarrow\space}}
}

\definecolor{XML_string}{RGB}{0,0,0}
\definecolor{XML_identifier}{RGB}{0,0,128}
\definecolor{XML_keyword}{RGB}{200,0,0}
\lstdefinelanguage{XML}
{
	morestring=[b]",
	morestring=[s]{>}{<},
	morecomment=[s]{<?}{?>},
	stringstyle=\color{XML_string},
	identifierstyle=\color{XML_identifier},
	keywordstyle=\color{XML_keyword},
	morekeywords={virtio}
}



% new command: header and footer
\newcommand{\virtioheaderfooter}{
\pagestyle{fancy}
\fancyhf{}
\renewcommand{\headrulewidth}{0pt}
\renewcommand{\footrulewidth}{0pt}
\fancyfoot[L]{\footnotesize\virtiospecfile\\\virtioworkproduct}
\fancyfoot[C]{\footnotesize\ \\\virtiooasiscopyright}
\fancyfoot[R]{\footnotesize\virtioworkingdraftdate\\Page \thepage\ of \pageref*{LastPage}}
}

% page setup
\setlength{\headheight}{0pt}
\setlength{\headsep}{10pt}
\setlength{\extrarowheight}{0.05in}

% headers setup
\titleformat{\chapter}[block]{\titlerule}{}{0pt}{\color{oasis1}\fontsize{18}{18}\textbf{\ifthenelse{\equal{\@chapapp}{\appendixname}}{Appendix~\thechapter.\ #1}{\ifthenelse{\equal{\thechapter}{0}}{}{\thechapter\ \ }#1}}\vspace{-0.3in}}
\titleformat{\section}{\color{oasis1}\normalfont\Large\bfseries}{\color{oasis1}\thesection}{1em}{#1}
\titleformat{\subsection}{\color{oasis1}\normalfont\large\bfseries}{\color{oasis1}\thesubsection}{1em}{#1}
\titleformat{\subsubsection}{\color{oasis1}\normalfont\normalsize\bfseries}{\color{oasis1}\thesubsubsection}{1em}{#1}
\titleformat{\paragraph}[hang]{\color{oasis1}\normalfont\normalsize\bfseries}{\color{oasis1}\theparagraph}{1em}{#1}
\titleformat{\subparagraph}[hang]{\color{oasis1}\normalfont\normalsize\bfseries}{\color{oasis1}\thesubparagraph}{1em}{#1}

% toc hacks
\makeatletter
% hyperref package: do not include the chapter/section number in the link
\def\process@contentsline#1#2{#1{#2}\hyper@linkstart{link}{\Hy@tocdestname}}
\patchcmd{\contentsline}{\hyper@linkstart{link}{\Hy@tocdestname}{#2}}{\process@contentsline#2}{}{}
% report style: add dotted lines between chapter names and page numbers
\patchcmd{\l@chapter}{\hfil}{\leaders\hbox{\normalfont$\m@th\mkern \@dotsep mu\hbox{.}\mkern \@dotsep mu$}\hfill}{}{}
\makeatother

% main document
% section numbers 5 levels deep
\setcounter{secnumdepth}{5}

\begin{document}
\pagestyle{fancy}

% start text
\setcounter{page}{1}

%%% define header and footer
\virtioheaderfooter

%%% define header an footer for chapter title pages
\fancypagestyle{plain}{\fancyhf{}\virtioheaderfooter}

% title page
\begin{titlepage}
\thispagestyle{fancy}

\begin{tabular*}{\textwidth}{@{\extracolsep{\fill}}lr}
	\includegraphics[height=0.55in]{images/oasis.png}\\
	\hline
\end{tabular*}

\oasistitle{\virtioversion}
\oasisstagelabel{\virtiodraftstagename\virtiodraftstageextratitle}
\vspace{0.05in}
\oasisstagelabel{\virtioworkingdraftdate}

\oasisspecificationuris{Specification URIs}

\begin{oasistitlesection}{This version}
\virtiourltex{ }(Authoritative)\newline
\virtiourlpdf\newline
\virtiourlhtml
\end{oasistitlesection}

\begin{oasistitlesection}{Previous version}
\url{http://docs.oasis-open.org/virtio/virtio/v1.0/cs01/tex/}
{}(Authoritative)\newline
\url{http://docs.oasis-open.org/virtio/virtio/v1.0/cs01/virtio-v1.0-cs01.pdf}\newline
\url{http://docs.oasis-open.org/virtio/virtio/v1.0/cs01/virtio-v1.0-cs01.html}
\end{oasistitlesection}

\begin{oasistitlesection}{Latest version}
\virtiourllatestpdf\newline
\virtiourllatesthtml
\end{oasistitlesection}

\begin{oasistitlesection}{Technical Committee}
\href{https://www.oasis-open.org/committees/virtio/}{OASIS Virtual I/O Device (VIRTIO) TC}
\end{oasistitlesection}

\begin{oasistitlesection}{Chairs}
Rusty Russell (\href{mailto:rusty@au.ibm.com}{rusty@au.ibm.com}), \href{http://www.ibm.com/}{IBM}\newline
Michael S. Tsirkin (\href{mailto:mst@redhat.com}{mst@redhat.com}), \href{http://www.redhat.com/}{Red Hat}\newline
\end{oasistitlesection}

\begin{oasistitlesection}{Editors}
Rusty Russell (\href{mailto:rusty@au.ibm.com}{rusty@au.ibm.com}), \href{http://www.ibm.com/}{IBM}\newline
Michael S. Tsirkin (\href{mailto:mst@redhat.com}{mst@redhat.com}), \href{http://www.redhat.com/}{Red Hat}\newline
Cornelia Huck (\href{mailto:cornelia.huck@de.ibm.com}{cornelia.huck@de.ibm.com}), \href{http://www.ibm.com/}{IBM}\newline
Pawel Moll (\href{mailto:pawel.moll@arm.com}{pawel.moll@arm.com}), \href{http://www.arm.com/}{ARM}
\end{oasistitlesection}


\begin{oasistitlesection}{Additional artifacts}
This prose specification is one component of a Work Product that also includes:
\begin{itemize*}
	\item Example Driver Listing: \newline
	\virtiourllistings
\end{itemize*}
\end{oasistitlesection}
\vspace{-0.2in}

\begin{oasistitlesection}{Related work}
This specification replaces or supersedes:
\begin{itemize*}
	\item Virtio PCI Card Specification Version 0.9.5:\newline
	\url{http://ozlabs.org/~rusty/virtio-spec/virtio-0.9.5.pdf}
\end{itemize*}
% This specification is related to:
% \begin{itemize*}
%        \item Related specifications (list)
% \end{itemize*}
\end{oasistitlesection}
\vspace{-0.2in}

% \begin{oasistitlesection}{Declared XML namespaces}
% \vspace{-0.1in}
% \begin{itemize*}
% 	\item namespaces which are declared, not just referenced or used (list)
% \end{itemize*}
% \end{oasistitlesection}
\vfill\newpage

\begin{oasistitlesection}{Abstract}
This document describes the specifications of the “virtio” family of
devices. These devices are found in virtual environments, yet by
design they look like physical devices to the guest
within the virtual machine - and this
document treats them as such. This similarity allows the guest to use standard
drivers and discovery mechanisms.

The purpose of virtio and this specification is that virtual
environments and guests should have a straightforward, efficient,
standard and extensible mechanism for virtual devices, rather
than boutique per-environment or per-OS mechanisms.

\end{oasistitlesection}

\begin{oasistitlesection}{Status}
This document was last revised or approved by the Virtual I/O Device
(VIRTIO) TC on the above date. The level of approval is also listed above. Check the ``Latest version'' location noted above for possible later revisions of this document.
Any other numbered Versions and other
technical work produced by the Technical Committee (TC) are
listed at
\url{https://www.oasis-open.org/committees/tc_home.php?wg_abbrev=virtio#technical}.

Technical Committee members should send comments on this
specification to the Technical Committee’s email list. Others
should send comments to the Technical Committee by using the
``\href{https://www.oasis-open.org/committees/comments/form.php?wg_abbrev=virtio}{Send
A Comment}'' button on the Technical Committee’s web page at
\url{https://www.oasis-open.org/committees/virtio/}.

For information on whether any patents have been disclosed that
may be essential to implementing this specification, and any
offers of patent licensing terms, please refer to the
Intellectual Property Rights section of the Technical Committee
web page (\url{https://www.oasis-open.org/committees/virtio/ipr.php}).
\end{oasistitlesection}


\begin{oasistitlesection}{Citation format}
When referencing this specification the following citation format should be used:\newline

\textbf{[VIRTIO-v\virtiorev]}\newline
\textit{\virtioversion}. Edited by Rusty Russell, Michael S.
Tsirkin, Cornelia Huck, and Pawel Moll. \virtioworkingdraftdate.
\virtiodraftoasisstagename \virtiodraftstageextra . \virtiourlhtml .
Latest version: \virtiourllatesthtml .
\end{oasistitlesection}

\vfill\newpage

\oasisnoticelabel{Notices}

Copyright © OASIS Open 2014. All Rights Reserved.

All capitalized terms in the following text have the meanings assigned
to them in the OASIS Intellectual Property Rights Policy (the "OASIS
IPR Policy"). The full \href{https://www.oasis-open.org/policies-guidelines/ipr}{Policy} may be found at the OASIS website.

This document and translations of it may be copied and furnished to
others, and derivative works that comment on or otherwise explain it
or assist in its implementation may be prepared, copied, published,
and distributed, in whole or in part, without restriction of any kind,
provided that the above copyright notice and this section are included
on all such copies and derivative works. However, this document itself
may not be modified in any way, including by removing the copyright
notice or references to OASIS, except as needed for the purpose of
developing any document or deliverable produced by an OASIS Technical
Committee (in which case the rules applicable to copyrights, as set
forth in the OASIS IPR Policy, must be followed) or as required to
translate it into languages other than English.

The limited permissions granted above are perpetual and will not be
revoked by OASIS or its successors or assigns.

This document and the information contained herein is provided on an
"AS IS" basis and OASIS DISCLAIMS ALL WARRANTIES, EXPRESS OR IMPLIED,
INCLUDING BUT NOT LIMITED TO ANY WARRANTY THAT THE USE OF THE
INFORMATION HEREIN WILL NOT INFRINGE ANY OWNERSHIP RIGHTS OR ANY
IMPLIED WARRANTIES OF MERCHANTABILITY OR FITNESS FOR A PARTICULAR
PURPOSE.

OASIS requests that any OASIS Party or any other party that believes
it has patent claims that would necessarily be infringed by
implementations of this OASIS Committee Specification or OASIS
Standard, to notify OASIS TC Administrator and provide an indication
of its willingness to grant patent licenses to such patent claims in a
manner consistent with the IPR Mode of the OASIS Technical Committee
that produced this specification.

OASIS invites any party to contact the OASIS TC Administrator if it is
aware of a claim of ownership of any patent claims that would
necessarily be infringed by implementations of this specification by a
patent holder that is not willing to provide a license to such patent
claims in a manner consistent with the IPR Mode of the OASIS Technical
Committee that produced this specification. OASIS may include such
claims on its website, but disclaims any obligation to do so.

OASIS takes no position regarding the validity or scope of any
intellectual property or other rights that might be claimed to pertain
to the implementation or use of the technology described in this
document or the extent to which any license under such rights might or
might not be available; neither does it represent that it has made any
effort to identify any such rights. Information on OASIS' procedures
with respect to rights in any document or deliverable produced by an
OASIS Technical Committee can be found on the OASIS website. Copies of
claims of rights made available for publication and any assurances of
licenses to be made available, or the result of an attempt made to
obtain a general license or permission for the use of such proprietary
rights by implementers or users of this OASIS Committee Specification
or OASIS Standard, can be obtained from the OASIS TC Administrator.
OASIS makes no representation that any information or list of
intellectual property rights will at any time be complete, or that any
claims in such list are, in fact, Essential Claims.

The name "OASIS" is a trademark of \href{https://www.oasis-open.org/}{OASIS}, the owner and developer of
this specification, and should be used only to refer to the
organization and its official outputs. OASIS welcomes reference to,
and implementation and use of, specifications, while reserving the
right to enforce its marks against misleading uses. Please see
\url{https://www.oasis-open.org/policies-guidelines/trademark} for above guidance.
\\\\

\end{titlepage}
 

% ToC
\setcounter{page}{4}
\setcounter{tocdepth}{4}
\renewcommand{\contentsname}{Table of Contents}
\addtocontents{toc}{\protect\virtioheaderfooter}
\tableofcontents
\newpage

%%% line numbers?
%\linenumbers

% introduction
\chapter{Introduction}

This document describes the specifications of the “virtio” family of
devices. These are devices are found in virtual environments, yet by
design they are not all that different from physical devices, and this
document treats them as such. This allows the guest to use standard
drivers and discovery mechanisms.

The purpose of virtio and this specification is that virtual
environments and guests should have a straightforward, efficient,
standard and extensible mechanism for virtual devices, rather
than boutique per-environment or per-OS mechanisms.

  Straightforward: Virtio devices use normal bus mechanisms of
  interrupts and DMA which should be familiar to any device driver
  author. There is no exotic page-flipping or COW mechanism: it's just
  a normal device.
\footnote{This lack of page-sharing implies that the implementation of the
device (e.g. the hypervisor or host) needs full access to the
guest memory. Communication with untrusted parties (i.e.
inter-guest communication) requires copying.
}

  Efficient: Virtio devices consist of rings of descriptors
  for input and output, which are neatly separated to avoid cache
  effects from both guest and device writing to the same cache
  lines.

  Standard: Virtio makes no assumptions about the environment in which
  it operates, beyond supporting the bus attaching the device.  Virtio
  devices are implemented over PCI and other buses, and earlier drafts
  been implemented on other buses not included in this spec.
\footnote{The Linux implementation further separates the PCI virtio code
from the specific virtio drivers: these drivers are shared with
the non-PCI implementations (currently lguest and S/390).
}

  Extensible: Virtio PCI devices contain feature bits which are
  acknowledged by the guest operating system during device setup.
  This allows forwards and backwards compatibility: the device
  offers all the features it knows about, and the driver
  acknowledges those it understands and wishes to use.

\section{Terminology}\label{Terminology}

The key words "MUST", "MUST NOT", "REQUIRED", "SHALL", "SHALL NOT", "SHOULD", "SHOULD NOT", "RECOMMENDED", "MAY", and "OPTIONAL" in this document are to be interpreted as described in \hyperref[intro:rfc2119]{[RFC2119]}.

\section{Normative References}

\begin{longtable}{l p{5in}}
	\phantomsection\label{intro:rfc2119}\textbf{[RFC2119]} & S. Bradner, Key words for use in RFCs to Indicate Requirement Levels, \newline\url{http://www.ietf.org/rfc/rfc2119.txt}, March 1997\\
	\phantomsection\label{intro:S390 PoP}\textbf{[S390 PoP]} & z/Architecture Principles of Operation, \newline IBM Publication SA22-7832\\
	\phantomsection\label{intro:S390 Common I/O}\textbf{[S390 Common I/O]} & ESA/390 Common I/O-Device and Self-Description, \newline IBM Publication SA22-7204\\
\end{longtable}


\section{Non-Normative References}


\newpage

 

% Conformance
%%%%%%%%%%%%%%%%%%%%%%%%%%%%%%%%%%%%%%%%%%%%%%%%%%%%%%%%%%%%
%% Conformance
%%%%%%%%%%%%%%%%%%%%%%%%%%%%%%%%%%%%%%%%%%%%%%%%%%%%%%%%%%%%
\chapter{Conformance}

This chapter lists the conformance targets and clauses for each; this
also forms a useful checklist which authors are asked to consult for their
implementations!

\section{Conformance Targets}\label{sec:Conformance / Conformance Targets}

Conformance targets:
\begin{description}
\item[Driver] A driver MUST conform to three conformance clauses:
  \begin{itemize}
    \item Clause \ref{sec:Conformance / Driver Conformance},
    \item One of clauses \ref{sec:Conformance / Driver Conformance / PCI Driver Conformance}, \ref{sec:Conformance / Driver Conformance / MMIO Driver Conformance} or \ref{sec:Conformance / Driver Conformance / Channel I/O Driver Conformance}.
    \item One of clauses \ref{sec:Conformance / Driver Conformance / Network Driver Conformance}, \ref{sec:Conformance / Driver Conformance / Block Driver Conformance}, \ref{sec:Conformance / Driver Conformance / Console Driver Conformance}, \ref{sec:Conformance / Driver Conformance / Entropy Driver Conformance}, \ref{sec:Conformance / Driver Conformance / Traditional Memory Balloon Driver Conformance} or \ref{sec:Conformance / Driver Conformance / SCSI Host Driver Conformance}.
  \end{itemize}
\item[Device] A device MUST conform to three conformance clauses:
  \begin{itemize}
    \item Clause \ref{sec:Conformance / Device Conformance},
    \item One of clauses \ref{sec:Conformance / Device Conformance / PCI Device Conformance}, \ref{sec:Conformance / Device Conformance / MMIO Device Conformance} or \ref{sec:Conformance / Device Conformance / Channel I/O Device Conformance}.
    \item One of clauses \ref{sec:Conformance / Device Conformance / Network Device Conformance}, \ref{sec:Conformance / Device Conformance / Block Device Conformance}, \ref{sec:Conformance / Device Conformance / Console Device Conformance}, \ref{sec:Conformance / Device Conformance / Entropy Device Conformance}, \ref{sec:Conformance / Device Conformance / Traditional Memory Balloon Device Conformance} or \ref{sec:Conformance / Device Conformance / SCSI Host Device Conformance}.
  \end{itemize}
\end{description}

\section{Driver Conformance}\label{sec:Conformance / Driver Conformance}

A driver MUST conform to the following normative statements:

\begin{itemize}
\item \ref{drivernormative:Basic Facilities of a Virtio Device / Device Status Field}
\item \ref{drivernormative:Basic Facilities of a Virtio Device / Feature Bits}
\item \ref{drivernormative:Basic Facilities of a Virtio Device / Device Configuration Space}
\item \ref{drivernormative:Basic Facilities of a Virtio Device / Virtqueues}
\item \ref{drivernormative:Basic Facilities of a Virtio Device / Message Framing}
\item \ref{drivernormative:Basic Facilities of a Virtio Device / Virtqueues / The Virtqueue Descriptor Table}
\item \ref{drivernormative:Basic Facilities of a Virtio Device / Virtqueues / The Virtqueue Descriptor Table / Indirect Descriptors}
\item \ref{drivernormative:Basic Facilities of a Virtio Device / Virtqueues / Used Buffer Notification Suppression}
\item \ref{drivernormative:Basic Facilities of a Virtio Device / Virtqueues / The Virtqueue Available Ring}
\item \ref{drivernormative:Basic Facilities of a Virtio Device / Virtqueues / The Virtqueue Used Ring}
\item \ref{drivernormative:Basic Facilities of a Virtio Device / Virtqueues / Available Buffer Notification Suppression}
\item \ref{drivernormative:Basic Facilities of a Virtio Device / Virtqueues / Supplying Buffers to The Device / Updating idx}
\item \ref{drivernormative:Basic Facilities of a Virtio Device / Virtqueues / Supplying Buffers to The Device / Notifying The Device}
\item \ref{drivernormative:General Initialization And Device Operation / Device Initialization}
\item \ref{drivernormative:General Initialization And Device Operation / Device Cleanup}
\item \ref{drivernormative:Reserved Feature Bits}
\end{itemize}

\subsection{PCI Driver Conformance}\label{sec:Conformance / Driver Conformance / PCI Driver Conformance}

A PCI driver MUST conform to the following normative statements:

\begin{itemize}
\item \ref{drivernormative:Virtio Transport Options / Virtio Over PCI Bus / PCI Device Discovery}
\item \ref{drivernormative:Virtio Transport Options / Virtio Over PCI Bus / PCI Device Layout}
\item \ref{drivernormative:Virtio Transport Options / Virtio Over PCI Bus / Virtio Structure PCI Capabilities}
\item \ref{drivernormative:Virtio Transport Options / Virtio Over PCI Bus / PCI Device Layout / Common configuration structure layout}
\item \ref{drivernormative:Virtio Transport Options / Virtio Over PCI Bus / PCI Device Layout / ISR status capability}
\item \ref{drivernormative:Virtio Transport Options / Virtio Over PCI Bus / PCI Device Layout / PCI configuration access capability}
\item \ref{drivernormative:Virtio Transport Options / Virtio Over PCI Bus / PCI-specific Initialization And Device Operation / Device Initialization / MSI-X Vector Configuration}
\item \ref{drivernormative:Virtio Transport Options / Virtio Over PCI Bus / PCI-specific Initialization And Device Operation / Notification of Device Configuration Changes}
\end{itemize}

\subsection{MMIO Driver Conformance}\label{sec:Conformance / Driver Conformance / MMIO Driver Conformance}

An MMIO driver MUST conform to the following normative statements:

\begin{itemize}
\item \ref{drivernormative:Virtio Transport Options / Virtio Over MMIO / MMIO Device Register Layout}
\item \ref{drivernormative:Virtio Transport Options / Virtio Over MMIO / MMIO-specific Initialization And Device Operation / Device Initialization}
\item \ref{drivernormative:Virtio Transport Options / Virtio Over MMIO / MMIO-specific Initialization And Device Operation / Notifications From The Device}
\end{itemize}

\subsection{Channel I/O Driver Conformance}\label{sec:Conformance / Driver Conformance / Channel I/O Driver Conformance}

A Channel I/O driver MUST conform to the following normative statements:

\begin{itemize}
\item \ref{drivernormative:Virtio Transport Options / Virtio over channel I/O / Basic Concepts}
\item \ref{drivernormative:Virtio Transport Options / Virtio over channel I/O / Device Initialization / Setting the Virtio Revision}
\item \ref{drivernormative:Virtio Transport Options / Virtio over channel I/O / Device Initialization / Communicating Status Information}
\item \ref{drivernormative:Virtio Transport Options / Virtio over channel I/O / Device Operation / Host->Guest Notification / Notification via Adapter I/O Interrupts}
\item \ref{drivernormative:Virtio Transport Options / Virtio over channel I/O / Device Operation / Guest->Host Notification}
\end{itemize}

\subsection{Network Driver Conformance}\label{sec:Conformance / Driver Conformance / Network Driver Conformance}

A network driver MUST conform to the following normative statements:

\begin{itemize}
\item \ref{drivernormative:Device Types / Network Device / Device configuration layout}
\item \ref{drivernormative:Device Types / Network Device / Device Operation / Packet Transmission}
\item \ref{drivernormative:Device Types / Network Device / Device Operation / Setting Up Receive Buffers}
\item \ref{drivernormative:Device Types / Network Device / Device Operation / Processing of Incoming Packets}
\item \ref{drivernormative:Device Types / Network Device / Device Operation / Control Virtqueue / Packet Receive Filtering}
\item \ref{drivernormative:Device Types / Network Device / Device Operation / Control Virtqueue / Setting MAC Address Filtering}
\item \ref{drivernormative:Device Types / Network Device / Device Operation / Control Virtqueue / Gratuitous Packet Sending}
\item \ref{drivernormative:Device Types / Network Device / Device Operation / Control Virtqueue / Automatic receive steering in multiqueue mode}
\item \ref{drivernormative:Device Types / Network Device / Device Operation / Control Virtqueue / Offloads State Configuration / Setting Offloads State}
\end{itemize}

\subsection{Block Driver Conformance}\label{sec:Conformance / Driver Conformance / Block Driver Conformance}

A block driver MUST conform to the following normative statements:

\begin{itemize}
\item \ref{drivernormative:Device Types / Block Device / Device Initialization}
\item \ref{drivernormative:Device Types / Block Device / Device Operation}
\end{itemize}

\subsection{Console Driver Conformance}\label{sec:Conformance / Driver Conformance / Console Driver Conformance}

A console driver MUST conform to the following normative statements:

\begin{itemize}
\item \ref{drivernormative:Device Types / Console Device / Device Operation}
\item \ref{drivernormative:Device Types / Console Device / Device Operation / Multiport Device Operation}
\end{itemize}

\subsection{Entropy Driver Conformance}\label{sec:Conformance / Driver Conformance / Entropy Driver Conformance}

An entropy driver MUST conform to the following normative statements:

\begin{itemize}
\item \ref{drivernormative:Device Types / Entropy Device / Device Operation}
\end{itemize}

\subsection{Traditional Memory Balloon Driver Conformance}\label{sec:Conformance / Driver Conformance / Traditional Memory Balloon Driver Conformance}

A traditional memory balloon driver MUST conform to the following normative statements:

\begin{itemize}
\item \ref{drivernormative:Device Types / Memory Balloon Device / Feature bits}
\item \ref{drivernormative:Device Types / Memory Balloon Device / Device Operation}
\item \ref{drivernormative:Device Types / Memory Balloon Device / Device Operation / Memory Statistics}
\end{itemize}

\subsection{SCSI Host Driver Conformance}\label{sec:Conformance / Driver Conformance / SCSI Host Driver Conformance}

An SCSI host driver MUST conform to the following normative statements:

\begin{itemize}
\item \ref{drivernormative:Device Types / SCSI Host Device / Device configuration layout}
\item \ref{drivernormative:Device Types / SCSI Host Device / Device Operation / Device Operation: Request Queues}
\item \ref{drivernormative:Device Types / SCSI Host Device / Device Operation / Device Operation: eventq}
\end{itemize}

\subsection{Input Driver Conformance}\label{sec:Conformance / Driver Conformance / Input Driver Conformance}

An input driver MUST conform to the following normative statements:

\begin{itemize}
\item \ref{drivernormative:Device Types / Input Device / Device Initialization}
\item \ref{drivernormative:Device Types / Input Device / Device Operation}
\end{itemize}

\subsection{Crypto Driver Conformance}\label{sec:Conformance / Driver Conformance / Crypto Driver Conformance}

A Crypto driver MUST conform to the following normative statements:

\begin{itemize}
\item \ref{drivernormative:Device Types / Crypto Device / Device configuration layout}
\item \ref{drivernormative:Device Types / Crypto Device / Device Initialization}
\item \ref{drivernormative:Device Types / Crypto Device / Device Operation / Control Virtqueue / Session operation / Session operation: create session}
\item \ref{drivernormative:Device Types / Crypto Device / Device Operation / Control Virtqueue / Session operation / Session operation: destroy session}
\item \ref{drivernormative:Device Types / Crypto Device / Device Operation / HASH Service Operation}
\item \ref{drivernormative:Device Types / Crypto Device / Device Operation / MAC Service Operation}
\item \ref{drivernormative:Device Types / Crypto Device / Device Operation / Symmetric algorithms Operation}
\item \ref{drivernormative:Device Types / Crypto Device / Device Operation / AEAD Service Operation}
\end{itemize}

\section{Device Conformance}\label{sec:Conformance / Device Conformance}

A device MUST conform to the following normative statements:

\begin{itemize}
\item \ref{devicenormative:Basic Facilities of a Virtio Device / Device Status Field}
\item \ref{devicenormative:Basic Facilities of a Virtio Device / Feature Bits}
\item \ref{devicenormative:Basic Facilities of a Virtio Device / Device Configuration Space}
\item \ref{devicenormative:Basic Facilities of a Virtio Device / Message Framing}
\item \ref{devicenormative:Basic Facilities of a Virtio Device / Virtqueues / The Virtqueue Descriptor Table}
\item \ref{devicenormative:Basic Facilities of a Virtio Device / Virtqueues / The Virtqueue Descriptor Table / Indirect Descriptors}
\item \ref{devicenormative:Basic Facilities of a Virtio Device / Virtqueues / Used Buffer Notification Suppression}
\item \ref{devicenormative:Basic Facilities of a Virtio Device / Virtqueues / The Virtqueue Used Ring}
\item \ref{devicenormative:Basic Facilities of a Virtio Device / Virtqueues / Available Buffer Notification Suppression}
\item \ref{devicenormative:Reserved Feature Bits}
\end{itemize}

\subsection{PCI Device Conformance}\label{sec:Conformance / Device Conformance / PCI Device Conformance}

A PCI device MUST conform to the following normative statements:

\begin{itemize}
\item \ref{devicenormative:Virtio Transport Options / Virtio Over PCI Bus}
\item \ref{devicenormative:Virtio Transport Options / Virtio Over PCI Bus / PCI Device Discovery}
\item \ref{devicenormative:Virtio Transport Options / Virtio Over PCI Bus / PCI Device Layout}
\item \ref{devicenormative:Virtio Transport Options / Virtio Over PCI Bus / Virtio Structure PCI Capabilities}
\item \ref{devicenormative:Virtio Transport Options / Virtio Over PCI Bus / PCI Device Layout / Common configuration structure layout}
\item \ref{devicenormative:Virtio Transport Options / Virtio Over PCI Bus / PCI Device Layout / Notification capability}
\item \ref{devicenormative:Virtio Transport Options / Virtio Over PCI Bus / PCI Device Layout / ISR status capability}
\item \ref{devicenormative:Virtio Transport Options / Virtio Over PCI Bus / PCI Device Layout / Device-specific configuration}
\item \ref{devicenormative:Virtio Transport Options / Virtio Over PCI Bus / PCI Device Layout / PCI configuration access capability}
\item \ref{devicenormative:Virtio Transport Options / Virtio Over PCI Bus / PCI-specific Initialization And Device Operation / Device Initialization / Non-transitional Device With Legacy Driver}
\item \ref{devicenormative:Virtio Transport Options / Virtio Over PCI Bus / PCI-specific Initialization And Device Operation / Device Initialization / MSI-X Vector Configuration}
\item \ref{devicenormative:Virtio Transport Options / Virtio Over PCI Bus / PCI-specific Initialization And Device Operation / Used Buffer Notifications}
\item \ref{devicenormative:Virtio Transport Options / Virtio Over PCI Bus / PCI-specific Initialization And Device Operation / Notification of Device Configuration Changes}
\end{itemize}

\subsection{MMIO Device Conformance}\label{sec:Conformance / Device Conformance / MMIO Device Conformance}

An MMIO device MUST conform to the following normative statements:

\begin{itemize}
\item \ref{devicenormative:Virtio Transport Options / Virtio Over MMIO / MMIO Device Register Layout}
\end{itemize}

\subsection{Channel I/O Device Conformance}\label{sec:Conformance / Device Conformance / Channel I/O Device Conformance}

A Channel I/O device MUST conform to the following normative statements:

\begin{itemize}
\item \ref{devicenormative:Virtio Transport Options / Virtio over channel I/O / Basic Concepts}
\item \ref{devicenormative:Virtio Transport Options / Virtio over channel I/O / Device Initialization / Setting the Virtio Revision}
\item \ref{devicenormative:Virtio Transport Options / Virtio over channel I/O / Device Initialization / Configuring a Virtqueue}
\item \ref{devicenormative:Virtio Transport Options / Virtio over channel I/O / Device Initialization / Communicating Status Information}
\item \ref{devicenormative:Virtio Transport Options / Virtio over channel I/O / Device Initialization / Setting Up Indicators / Setting Up Two-Stage Queue Indicators}
\item \ref{devicenormative:Virtio Transport Options / Virtio over channel I/O / Device Operation / Host->Guest Notification / Notification via Adapter I/O Interrupts}
\item \ref{devicenormative:Virtio Transport Options / Virtio over channel I/O / Device Operation / Guest->Host Notification}
\end{itemize}

\subsection{Network Device Conformance}\label{sec:Conformance / Device Conformance / Network Device Conformance}

A network device MUST conform to the following normative statements:

\begin{itemize}
\item \ref{devicenormative:Device Types / Network Device / Device configuration layout}
\item \ref{devicenormative:Device Types / Network Device / Device Operation / Packet Transmission}
\item \ref{devicenormative:Device Types / Network Device / Device Operation / Setting Up Receive Buffers}
\item \ref{devicenormative:Device Types / Network Device / Device Operation / Processing of Incoming Packets}
\item \ref{devicenormative:Device Types / Network Device / Device Operation / Control Virtqueue / Packet Receive Filtering}
\item \ref{devicenormative:Device Types / Network Device / Device Operation / Control Virtqueue / Setting MAC Address Filtering}
\item \ref{devicenormative:Device Types / Network Device / Device Operation / Control Virtqueue / Gratuitous Packet Sending}
\item \ref{devicenormative:Device Types / Network Device / Device Operation / Control Virtqueue / Automatic receive steering in multiqueue mode}
\end{itemize}

\subsection{Block Device Conformance}\label{sec:Conformance / Device Conformance / Block Device Conformance}

A block device MUST conform to the following normative statements:

\begin{itemize}
\item \ref{devicenormative:Device Types / Block Device / Device Initialization}
\item \ref{devicenormative:Device Types / Block Device / Device Operation}
\end{itemize}

\subsection{Console Device Conformance}\label{sec:Conformance / Device Conformance / Console Device Conformance}

A console device MUST conform to the following normative statements:

\begin{itemize}
\item \ref{devicenormative:Device Types / Console Device / Device Initialization}
\item \ref{devicenormative:Device Types / Console Device / Device Operation / Multiport Device Operation}
\end{itemize}

\subsection{Entropy Device Conformance}\label{sec:Conformance / Device Conformance / Entropy Device Conformance}

An entropy device MUST conform to the following normative statements:

\begin{itemize}
\item \ref{devicenormative:Device Types / Entropy Device / Device Operation}
\end{itemize}

\subsection{Traditional Memory Balloon Device Conformance}\label{sec:Conformance / Device Conformance / Traditional Memory Balloon Device Conformance}

A traditional memory balloon device MUST conform to the following normative statements:

\begin{itemize}
\item \ref{devicenormative:Device Types / Memory Balloon Device / Feature bits}
\item \ref{devicenormative:Device Types / Memory Balloon Device / Device Operation}
\item \ref{devicenormative:Device Types / Memory Balloon Device / Device Operation / Memory Statistics}
\end{itemize}

\subsection{SCSI Host Device Conformance}\label{sec:Conformance / Device Conformance / SCSI Host Device Conformance}

An SCSI host device MUST conform to the following normative statements:

\begin{itemize}
\item \ref{devicenormative:Device Types / SCSI Host Device / Device configuration layout}
\item \ref{devicenormative:Device Types / SCSI Host Device / Device Initialization}
\item \ref{devicenormative:Device Types / SCSI Host Device / Device Operation / Device Operation: Request Queues}
\item \ref{devicenormative:Device Types / SCSI Host Device / Device Operation / Device Operation: eventq}
\end{itemize}

\subsection{Input Device Conformance}\label{sec:Conformance / Device Conformance / Input Device Conformance}

An input device MUST conform to the following normative statements:

\begin{itemize}
\item \ref{devicenormative:Device Types / Input Device / Device Initialization}
\item \ref{devicenormative:Device Types / Input Device / Device Operation}
\end{itemize}

\subsection{Crypto Device Conformance}\label{sec:Conformance / Device Conformance / Crypto Device Conformance}

A Crypto device MUST conform to the following normative statements:

\begin{itemize}
\item \ref{devicenormative:Device Types / Crypto Device / Device configuration layout}
\item \ref{devicenormative:Device Types / Crypto Device / Device Operation / Control Virtqueue / Session operation / Session operation: create session}
\item \ref{devicenormative:Device Types / Crypto Device / Device Operation / Control Virtqueue / Session operation / Session operation: destroy session}
\item \ref{devicenormative:Device Types / Crypto Device / Device Operation / HASH Service Operation}
\item \ref{devicenormative:Device Types / Crypto Device / Device Operation / MAC Service Operation}
\item \ref{devicenormative:Device Types / Crypto Device / Device Operation / Symmetric algorithms Operation}
\item \ref{devicenormative:Device Types / Crypto Device / Device Operation / AEAD Service Operation}
\end{itemize}

\section{Legacy Interface: Transitional Device and
Transitional Driver Conformance}\label{sec:Conformance / Legacy
Interface: Transitional Device and 
Transitional Driver Conformance}
A conformant implementation MUST be either transitional or
non-transitional, see \ref{intro:Legacy
Interface: Terminology}.

A non-transitional implementation conforms to this specification
if it satisfies all of the MUST or REQUIRED level requirements
defined above.

An implementation MAY choose to implement OPTIONAL support for the
legacy interface, including support for legacy drivers
or devices, by additionally conforming to all of the MUST or
REQUIRED level requirements for the legacy interface
for the transitional devices and drivers.

The requirements for the legacy interface for transitional implementations
are located in sections named ``Legacy Interface'' listed below:
\begin{itemize}
\item Section \ref{sec:Basic Facilities of a Virtio Device /
Feature Bits / Legacy Interface: A Note on Feature Bits}
\item Section \ref{sec:Basic Facilities of a Virtio Device / Device Configuration Space / Legacy Interface: A Note on Configuration Space endian-ness}
\item Section \ref{sec:Basic Facilities of a Virtio Device / Device Configuration Space / Legacy Interface: Device Configuration Space}
\item Section \ref{sec:Basic Facilities of a Virtio Device / Virtqueues / Legacy Interfaces: A Note on Virtqueue Layout}
\item Section \ref{sec:Basic Facilities of a Virtio Device / Virtqueues / Legacy Interfaces: A Note on Virtqueue Endianness}
\item Section \ref{sec:Basic Facilities of a Virtio Device / Virtqueues / Message Framing / Legacy Interface: Message Framing}
\item Section \ref{sec:General Initialization And Device Operation / Device Initialization / Legacy Interface: Device Initialization}
\item Section \ref{sec:Virtio Transport Options / Virtio Over PCI Bus / PCI Device Discovery / Legacy Interfaces: A Note on PCI Device Discovery}
\item Section \ref{sec:Virtio Transport Options / Virtio Over PCI Bus / PCI Device Layout / Legacy Interfaces: A Note on PCI Device Layout}
\item Section \ref{sec:Virtio Transport Options / Virtio Over PCI Bus / PCI-specific Initialization And Device Operation / Device Initialization / Virtio Device Configuration Layout Detection / Legacy Interface: A Note on Device Layout Detection}
\item Section \ref{sec:Virtio Transport Options / Virtio Over PCI Bus / PCI-specific Initialization And Device Operation / Device Initialization / Virtqueue Configuration / Legacy Interface: A Note on Virtqueue Configuration}
\item Section \ref{sec:Virtio Transport Options / Virtio Over MMIO / Legacy interface}
\item Section \ref{sec:Virtio Transport Options / Virtio over channel I/O / Device Initialization / Setting the Virtio Revision / Legacy Interfaces: A Note on Setting the Virtio Revision}
\item Section \ref{sec:Virtio Transport Options / Virtio over channel I/O / Device Initialization / Configuring a Virtqueue / Legacy Interface: A Note on Configuring a Virtqueue}
\item Section \ref{sec:Virtio Transport Options / Virtio over channel I/O / Device Operation / Host->Guest Notification / Legacy Interfaces: A Note on Host->Guest Notification}
\item Section \ref{sec:Virtio Transport Options / Virtio over channel I/O / Device Initialization / Setting Up Indicators / Legacy Interfaces: A Note on Setting Up Indicators}
\item Section \ref{sec:Device Types / Network Device / Feature bits / Legacy Interface: Feature bits}
\item Section \ref{sec:Device Types / Network Device / Device configuration layout / Legacy Interface: Device configuration layout}
\item Section \ref{sec:Device Types / Network Device / Device Operation / Legacy Interface: Device Operation}
\item Section \ref{sec:Device Types / Network Device / Device Operation / Control Virtqueue / Setting MAC Address Filtering / Legacy Interface: Setting MAC Address Filtering}
\item Section \ref{sec:Device Types / Network Device / Device Operation / Control Virtqueue / VLAN Filtering / Legacy Interface: VLAN Filtering}
\item Section \ref{sec:Device Types / Network Device / Device Operation / Control Virtqueue / Automatic receive steering in multiqueue mode / Legacy Interface: Automatic receive steering in multiqueue mode}
\item Section \ref{sec:Device Types / Network Device / Device Operation / Control Virtqueue / Offloads State Configuration / Setting Offloads State / Legacy Interface: Setting Offloads State}
\item Section \ref{sec:Device Types / Block Device / Feature bits / Legacy Interface: Feature bits}
\item Section \ref{sec:Device Types / Block Device / Device configuration layout / Legacy Interface: Device configuration layout}
\item Section \ref{sec:Device Types / Block Device / Device Initialization / Legacy Interface: Device Initialization}
\item Section \ref{sec:Device Types / Block Device / Device Operation / Legacy Interface: Device Operation}
\item Section \ref{sec:Device Types / Console Device / Device configuration layout / Legacy Interface: Device configuration layout}
\item Section \ref{sec:Device Types / Console Device / Device Operation / Legacy Interface: Device Operation}
\item Section \ref{sec:Device Types / Memory Balloon Device / Feature bits / Legacy Interface: Feature bits}
\item Section \ref{sec:Device Types / Memory Balloon Device / Device Operation / Legacy Interface: Device Operation}
\item Section \ref{sec:Device Types / Memory Balloon Device / Device Operation / Memory Statistics / Legacy Interface: Memory Statistics}
\item Section \ref{sec:Device Types / SCSI Host Device / Device configuration layout / Legacy Interface: Device configuration layout}
\item Section \ref{sec:Device Types / SCSI Host Device / Device Operation / Legacy Interface: Device Operation}
\item Section \ref{sec:Device Types / SCSI Host Device / Device Operation / Device Operation: Request Queues / Legacy Interface: Device Operation: Request Queues}
\item Section \ref{sec:Device Types / SCSI Host Device / Device Operation / Device Operation: controlq / Legacy Interface: Device Operation: controlq}
\item Section \ref{sec:Device Types / SCSI Host Device / Device Operation / Device Operation: eventq / Legacy Interface: Device Operation: eventq}
\item Section \ref{sec:Reserved Feature Bits / Legacy Interface: Reserved Feature Bits}
\end{itemize}


% appendix
\appendix
\def\@chapapp{\appendixname}

% acknowledgements
%%%%%%%%%%%%%%%%%%%%%%%%%%%%%%%%%%%%%%%%%%%%%%%%%%%%%%%%%%%%
%% Acknowledgements
%%%%%%%%%%%%%%%%%%%%%%%%%%%%%%%%%%%%%%%%%%%%%%%%%%%%%%%%%%%%
\chapter{Acknowledgements}

The following individuals have participated in the creation of this specification and are gratefully acknowledged:

\begin{oasistitlesection}{Participants}
Amit Shah,	Red Hat	\newline
Amos Kong,	Red Hat	\newline
Anthony Liguori,	IBM	\newline
Bruce Rogers,	Novell	\newline
Bryan Venteicher,	NetApp	\newline
Cornelia Huck,	IBM	\newline
Daniel Kiper,	Oracle	\newline
Geoff Brown,	Machine-to-Machine Intelligence (M2MI) Corporation	\newline
Gershon Janssen,	Individual Member	\newline
James Bottomley,	Parallels IP Holdings GmbH	\newline
Luiz Capitulino,	Red Hat	\newline
Michael Tsirkin,	Red Hat	\newline
Paolo Bonzini,	Red Hat	\newline
Pawel Moll,	ARM Limited	\newline
Rusty Russell,	IBM	\newline
Sasha Levin,	Oracle	\newline
Sergey Tverdyshev,	Thales e-Security	\newline
Stefan Hajnoczi,	Red Hat	\newline
Tom Lyon,	Samya Systems, Inc.	\newline
\end{oasistitlesection}
 

%appendix stub
\chapter{Example Title}
demonstrating appendix numbering
\section{subsiduary appendix}
text
\subsection{Sub-subsiduary appendix}
text


% draft change log
\chapter{Revision History}
The following changes have been made since the previous version
of this specification:

\begin{longtable}{ | c | c | c | p{0.4\textwidth} | }
\hline
%\textbf{Revision} & \textbf{Date} & \textbf{Editor} & \textbf{Changes Made} \\
\hline
\endhead
%\lbrack Rev number \rbrack & \lbrack Rev Date \rbrack & \lbrack Modified By \rbrack & \lbrack Summary of Changes \rbrack \\
%\hline
478 & 15 Mar 2015 & Cornelia Huck & {VIRTIO-129: legacy:
clean up virtqueue layout definitions

Generalize "Legacy Interfaces: A Note on Virtqueue Layout" to allow
for different alignment requirements. Have pci and ccw refer to that
section for legacy devices. Remove the double definition of virtqueue
alignment (which referred to legacy, but was not tagged as such) from
the ccw section.
See \ref{sec:Basic Facilities of a Virtio Device / Virtqueues /
Legacy Interfaces: A Note on Virtqueue Layout}, \ref{sec:Virtio
Transport Options / Virtio Over PCI Bus / PCI-specific
Initialization And Device Operation / Device Initialization /
Virtqueue Configuration / Legacy Interface: A Note on Virtqueue
Configuration} and \ref{sec:Virtio Transport Options / Virtio
over channel I/O / Device Initialization / Configuring a
Virtqueue / Legacy Interface: A Note on Configuring a Virtqueue}.
 } \\
\hline
479 & 15 Mar 2015 & Cornelia Huck & {VIRTIO-116:
ccw: allow WRITE_STATUS to fail
    
We want to be able to fail setting a status on the device
(e.g.  FEATURES_OK if the device can't work with the features
negotiated).
The easiest way to do that is to allow the device to fail the
WRITE_STATUS command by posting a command reject.
See \ref{sec:Virtio Transport Options / Virtio over channel I/O /
Device Initialization / Communicating Status Information}.
 } \\
\hline
485 & 15 Mar 2015 & Jason Wang & {VIRTIO-135:
virtio-ring: comment fixup
    
virtio_ring.h included with spec has this text:
/* Support for avail_idx and used_idx fields */
it should really refer to avail_event and used_event.
See \ref{sec:virtio-ring.h}.
 } \\
\hline
486 & 15 Mar 2015 & Jason Wang & {VIRTIO-136:
document idx field in virtqueue used ring

Section \ref{sec:Basic Facilities of a Virtio Device / Virtqueues
/ The Virtqueue Used Ring} The Virtqueue Used Ring
listed the idx field, but never documented it.
See \ref{sec:Basic Facilities of a Virtio Device / Virtqueues /
The Virtqueue Used Ring}.
 } \\
\hline

\end{longtable}
 

\end{document} 

